\subsection{Relation to the instanton prediction}\label{subsec:instanton-model}
The denominator \(H\) in \eqref{eq:C4} is the normalized ordinary Weil--Petersson instanton action. To see the correspondence, choose the positive sign convention
\[
 A_{\rm inst}(w)=\frac{\sin(2\pi w)-2\pi w\cos(2\pi w)}{4\pi^3}
     =\frac23 w^3 H(2\pi^2w^2).
\]
With $x=w^2$, this is the effective potential of~\cite[equation~(3.15)]{EGGLS}, whose values at its stationary points are the instanton actions. The ordinary-volume polynomial prediction in~\cite[Section~3.3]{GPRS} and the one-boundary term in~\cite[equation~(4.22)]{JR} contain, up to an overall factor independent of the coefficient degree, the contour expression
\[
 \frac{\Gamma(h)}{L}\oint
       \frac{\mathrm{e}^{Lw}}{w A_{\rm inst}(w)^h}\,\mathrm{d}w,\qquad h=2g-1.
\]
The coefficient of \(L^{2k}\), multiplied by \((2k+1)!\), is proportional to
\[
 \Gamma(h)\operatorname*{Res}_{w=0} w^{2k}A_{\rm inst}(w)^{-h}\,\mathrm{d}w.
\]
Since \(3h-2k-1=2m\), expansion at zero gives the exact identity
\[
 \Gamma(h)\operatorname*{Res}_{w=0} w^{2k}A_{\rm inst}(w)^{-h}\,\mathrm{d}w
 =\Gamma(h)(3/2)^h(2\pi^2)^m \widetilde A_m(h)
 =2\,4^{-k}M_{g,k}.
\]
Thus the model for the coefficients \(a_{g,k}\) has precisely the normalization \(4^{-k}M_{g,k}\), up to the degree-independent overall factor in the prediction. Moreover, the exponent in the \(\mathrm{d}w/w\) representation has its stationary point at
\[
 \frac{2k+1}{h}=\frac{wA_{\rm inst}'(w)}{A_{\rm inst}(w)}
 =\frac{t^2\sin t}{\sin t-t\cos t},\qquad t=2\pi w.
\]
Thus both the comparison model and the saddle equation already appear in the instanton calculus. The estimates \eqref{eq:C5} and the concavity argument above justify the asymptotic formula for adjacent ratios of the true mixed intersection numbers, including both endpoint regimes; that proof does not use the contour expression.
